\chapter{Переоткрытие через нелинейный дискретный родитель}
\label{ch:v6-spectral-transition-discrete-nonlinear-parent-reopening-gate}

\section{Минимальный объект проверки}

Ретроспективная коррекция оставила один действительно непостроенный класс:
локальный дискретный закон, который не является сеточной записью заранее
заданного потенциала. Рассмотрим одномерный минимальный носитель
\begin{equation}
 \Psi_n^t\in\mathbb C^2_{\rm dir}\otimes E,
 \qquad E=M_{20\times15}(\mathbb C),
 \qquad \dim_{\mathbb C}E=300.
\end{equation}
Две компоненты отвечают встречным направлениям переноса, а $E$ сохраняет
полный внутренний носитель проекта.

Линейные локальные унитарные автоматы могут иметь уравнение Дирака как
длинноволновый предел \cite{discrete-parent-bialynicki1994,
discrete-parent-bisio2015}. Нелинейные квантовые блуждания со
state-dependent монетой аналогично дают нелинейный дираковский предел
\cite{discrete-parent-angles2023,discrete-parent-continuous-limits2019}.

\section{Точное нелинейное обновление}

Введём локальный вещественный билинеар
\begin{equation}
 s_n(\Psi)=
 \langle\Psi_n,(\sigma_z\otimes I_E)\Psi_n\rangle
\end{equation}
и монету
\begin{equation}
 C_{\kappa,n}(\Psi)=
 \exp\!\left[-i\bigl(\theta_0+\kappa s_n(\Psi)\bigr)
 \sigma_y\right]\otimes I_E.
 \label{eq:v6-discrete-parent-coin}
\end{equation}
Один шаг равен
\begin{equation}
 \Psi^{t+1}=S\,C_\kappa(\Psi^t)\Psi^t,
 \label{eq:v6-discrete-parent-update}
\end{equation}
где $S$ сдвигает правую компоненту на один узел вправо, а левую --- на
один узел влево.

Для каждого состояния $s_n$ вещественно, поэтому монета унитарна. Сдвиг
является перестановкой узлов. Следовательно,
\begin{equation}
 \sum_n\|\Psi_n^{t+1}\|^2=\sum_n\|\Psi_n^t\|^2
\end{equation}
точно, несмотря на нелинейность. Правило использует только один узел и его
ближайшие выходы, то есть проходит тесты локальности и сохранения нормы.

\section{Непрерывный предел}

При шаге решётки и времени $a$, масштабировании
\begin{equation}
 \theta_0=ma,\qquad \kappa=ga
\end{equation}
формальное разложение \eqref{eq:v6-discrete-parent-update} даёт оператор
типа
\begin{equation}
 i\partial_t\Psi=
 \left[-i\sigma_z\partial_x+m\sigma_y
 +g\langle\Psi,\sigma_z\Psi\rangle\sigma_y\right]\Psi+O(a).
 \label{eq:v6-discrete-parent-nonlinear-dirac}
\end{equation}
Тем самым дискретный закон действительно объединяет свободный дираковский
перенос и локальную самосвязь без отдельного непрерывного поля $q(x)$.

Однако параметры $m$, $g$ и физический шаг $a$ не следуют из унитарности,
локальности и однородности. Разные $\kappa$ дают разные режимы
самофокусировки; литература также рассматривает силу нелинейности как
параметр монеты, а не следствие одной только дискретности.

\section{Точный внутренний запрет}

На полном факторном соответствии уже доказано
\begin{equation}
 \operatorname{End}_{M_{20}-M_{15}}(E)=\mathbb C I_E.
 \label{eq:v6-discrete-parent-full-commutant}
\end{equation}
Поэтому всякая монета, совместимая с обоими полными действиями, имеет
внутреннюю часть, пропорциональную $I_E$. Более того, правило
\eqref{eq:v6-discrete-parent-update} ковариантно относительно любого
$V\in U(E)$:
\begin{equation}
 F(\Psi V)=F(\Psi)V.
\end{equation}
Оно не различает внутренние направления. Локализованный режим, если он
возникает в одном канале, существует в 300 эквивалентных копиях и их
унитарных суперпозициях.

Нелинейная зависимость от полной плотности может связать амплитуды каналов,
но не выбирает каноническую линию. Для сокращения кратности требуется
нескалярная монета, например проекторного типа. На полном бимодуле она
запрещена \eqref{eq:v6-discrete-parent-full-commutant}; после ограничения
до физической алгебры возвращается уже проверенная цепочка
$P_0\otimes P_\nu(H)$.

\begin{theorem}[Минимальный дискретный родитель не выбирает частицу]
Локальное state-dependent унитарное правило
\eqref{eq:v6-discrete-parent-update} точно сохраняет норму и имеет
нелинейный дираковский непрерывный предел. Но совместимость с полным
$M_{20}$--$M_{15}$-бимодулем делает внутреннюю монету скалярной. Поэтому
правило сохраняет кратность $300$, не фиксирует $g$ и $a$ и не выводит
единственный физический endpoint. Оно является новым дискретным
proof-of-concept, но не проходит R4 и R5.
\end{theorem}

\begin{nogocriterion}
Нельзя считать минимальный шаг решётки предсказанным размером частицы,
пока не выведена физическая величина $a$. Нельзя также вставлять
$P_0\otimes P_\nu(H)$ в монету без вывода из локальной алгебры обновления:
это воспроизводит старый операторнозначный дефект, а не решает его.
\end{nogocriterion}

\section{Следующий гейт}

Следующий гейт
\path{version6_spectral_transition_discrete_equivariant_coin_selector_gate}
должен классифицировать локальные монеты после перехода от полного фактора
к физической алгебре Стандартной модели и проверить, возникает ли
единственный нескалярный селектор без ручного выбора блока и Хиггса.

Нулевая гипотеза: коммутант физической алгебры даёт несколько центральных
блоков и не выбирает один из них. Стоп-критерий: если ранг-один монета
требует состояния Хиггса или внешнего проектора, дискретная ветвь
возвращается к условной лестнице Тома V и закрывается.

\begin{protocol}
\begin{enumerate}
 \item точное локальное нелинейное правило: \textbf{построено};
 \item сохранение нормы: \textbf{точное};
 \item ближайшая локальность: \textbf{да};
 \item нелинейный дираковский предел: \textbf{получен формально};
 \item минимальный дискретный шаг: \textbf{имеется};
 \item физическое значение шага: \textbf{не выведено};
 \item нелинейная связь $\kappa$: \textbf{не фиксирована};
 \item полная внутренняя ковариантность: \textbf{сохранена};
 \item внутренняя кратность: \textbf{$300$};
 \item единственный физический канал: \textbf{не выбран};
 \item статус: \textbf{дискретный proof-of-concept без физического замыкания}.
\end{enumerate}
\end{protocol}

Машинный сертификат сохранён в
\path{s2t_v6_spectral_transition_discrete_nonlinear_parent_reopening_gate_results.json}.

\begin{thebibliography}{9}
\bibitem{discrete-parent-bialynicki1994}
I.~Bialynicki-Birula,
\emph{Weyl, Dirac, and Maxwell Equations on a Lattice as Quantum Cellular Automata},
Physical Review D 49 (1994), 6920--6927; arXiv:hep-th/9304070.

\bibitem{discrete-parent-bisio2015}
A.~Bisio, G.~M.~D'Ariano, A.~Tosini,
\emph{Quantum Field as a Quantum Cellular Automaton: The Dirac Free Evolution in One Dimension},
Annals of Physics 354 (2015), 244--264; arXiv:1212.2839.

\bibitem{discrete-parent-angles2023}
A.~Angl\`es-Castillo, A.~P\'erez, E.~Rold\'an,
\emph{Solitons in a Photonic Nonlinear Quantum Walk: Lessons from the Continuum},
arXiv:2308.01014.

\bibitem{discrete-parent-continuous-limits2019}
M.~Maeda, A.~Suzuki,
\emph{Continuous Limits of Linear and Nonlinear Quantum Walks},
Reviews in Mathematical Physics 32 (2020), 2050008;
arXiv:1902.02017.
\end{thebibliography}